% Original Markdown SHA-256: 1954d499b7c63eb3f0202366ea75ff7f63729e17be110e11e7c6e3e9b425f650
\chapter{Exact composition, finite carry certificates, and the general-k capacity}
\label{chapter:03_general-k-carry-capacity}
{\small\noindent\textbf{Source:} \nolinkurl{history/EP817_GENERAL_K_POLYCLANK_20260914/research/general_k/notes/general-k-carry-capacity.md}\
\textbf{Identity:} \nolinkurl{1954d499b7c63eb3f0202366ea75ff7f63729e17be110e11e7c6e3e9b425f650}}
\medskip
\textbf{Erdős Problem 817 --- PolyClank research contribution, 14 September 2026.}

Base revision: \nolinkurl{23c0110c95b5a2036bdc04a1f352b6e5e27742c8} of \nolinkurl{KokunoYumeto/erdos-problem-817-workbench}.

Status: complete ordinary proofs of the statements below, with reproducible integer checks and explicit finite certificates; independent review is requested. Lean coverage remains that of the existing Core/Extended modules. This work adds no Lean declarations or replacement receipts. The conclusions below use no conjectural block-classification lemma.

\section{1. Scope and exact objects}\label{n03-scope-and-exact-objects}

For a finite set A of distinct positive integers, put

\[
\begin{adjustbox}{max width=.92\linewidth}
$\displaystyle H(A)=\{\sum_{a\in U}a:U\subseteq A\},\quad
S(A)=\sum_{a\in A}a,\quad R(A)=2S(A)+1.$
\end{adjustbox}
\]

Let \(\mathcal A_k(n)\) be the set of such A with \textbar A\textbar=n and with H(A) containing no nonconstant k-term arithmetic progression. Here k\textgreater=3. Set \(\mathcal A_k(0)=\{\varnothing\}\), and define

\[
\begin{adjustbox}{max width=.92\linewidth}
$\displaystyle g_k(n)=\min_{A\in\mathcal A_k(n)}\max A\quad(n\ge1),\qquad
h_k(n)=\min_{A\in\mathcal A_k(n)}R(A)\quad(n\ge0).$
\end{adjustbox}
\]

For a subset X of an abelian group, \(\operatorname{AP}_k(X)\) denotes the set of ordered tuples \((x_0,\ldots,x_{k-1})\in X^k\) with \(x_{i+2}-2x_{i+1}+x_i=0\) for every i. Constant tuples are included. Over a cyclic group, a nonconstant progression means that its common difference is a nonzero group element; its entries may repeat when that element has small order. All such orders are retained in the computations.

These definitions preserve every generator, its actual value, and every subset representation. R is an additional integer-valued observable. No gcd is removed, no generator is rescaled to 1, and no equality of subset sums is discarded. The original k=4 construction and signed-block proof remain credited to the anonymous/deleted Reddit contributor. The Lean finite upper bound remains separately credited to sneed-and-feed. Draft PR \#2 supplies the earlier variance investigation; the present proofs have no dependency on its new lemma.

The symbol lambda\_k below denotes the exponential growth rate. Korsky\textquotesingle s Equation (1.5) uses the same letter for the polynomial-loss exponent log\_2((k-1)/(k-2)); that quantity is kept separately here. The map log\_2:(1,infinity)-\textgreater(0,infinity), with inverse x-\textgreater2\^{}x, relates his polynomial exponent to the lower exponential base d\_k=(k-1)/(k-2).

The broader problem requests estimates for all k. Costa (2026) establishes \(\liminf g_3(n)/3^n=0\), addressing the historical constant-factor question. The present target is the full family of exponential rates and constructions.

\section{2. Exact product on generators and arithmetic-progression tuples}\label{n03-exact-product-on-generators-and-arithmetic-progression-tuples}

\textbf{Theorem 2.1 (composition with explicit inverse maps).} For arbitrary finite sets A,B of distinct positive integers, including the empty set, define

\[
\begin{adjustbox}{max width=.92\linewidth}
$\displaystyle A\star B=A\cup R(A)B.$
\end{adjustbox}
\]

Then

\[
\begin{adjustbox}{max width=.92\linewidth}
$\displaystyle |A\star B|=|A|+|B|,\qquad
S(A\star B)=S(A)+R(A)S(B),\qquad
R(A\star B)=R(A)R(B).$
\end{adjustbox}
\]

The maps

\[
\begin{adjustbox}{max width=.92\linewidth}
$\displaystyle \begin{aligned}
F_{A,B}:\mathcal P(A)\times\mathcal P(B)&\longrightarrow\mathcal P(A\star B),
 &(U,V)&\longmapsto U\cup R(A)V,\\
E_{A,B}:H(A)\times H(B)&\longrightarrow H(A\star B),
 &(x,y)&\longmapsto x+R(A)y
\end{aligned}$
\end{adjustbox}
\]

are bijections. E induces a bijection

\[
\begin{adjustbox}{max width=.92\linewidth}
$\displaystyle \operatorname{AP}_k(H(A))\times\operatorname{AP}_k(H(B))
\longrightarrow \operatorname{AP}_k(H(A\star B))$
\end{adjustbox}
\]

by componentwise evaluation. A target progression is constant exactly when both source progressions are constant. Consequently star sends \(\mathcal A_k(n)\times\mathcal A_k(m)\) into \(\mathcal A_k(n+m)\). It is associative and has identity \(\varnothing\).

\textbf{Proof.} Put R=R(A). Every member of A is less than R, while each member of RB is at least R. The two parts are disjoint, and multiplication by R is injective. The inverse of F sends W to \((W\cap A,\{b\in B:Rb\in W\})\). This proves the cardinality formula and the subset representation identity.

For x in H(A), 0\textless=x\textless=S(A)\textless R. Thus the inverse of E is the restriction of the Euclidean division map

\[
\begin{adjustbox}{max width=.92\linewidth}
$\displaystyle z\longmapsto \left(z-R\lfloor z/R\rfloor,\lfloor z/R\rfloor\right).$
\end{adjustbox}
\]

The first and second coordinates belong to H(A),H(B) because z arose from a subset of A star B. This proves the bijection on actual values. The subset-sum maps commute with F and E. In particular, writing mu\_A(x) for the number of subsets of A summing to x, one has the exact fiber identity \(\mu_{A\star B}(x+Ry)=\mu_A(x)\mu_B(y)\).

Take an arithmetic progression z\_i=x\_i+Ry\_i in the target and put \(e_i=x_{i+2}-2x_{i+1}+x_i\). The equation for z gives \(e_i=-R(y_{i+2}-2y_{i+1}+y_i)\). Meanwhile \(-2S(A)\le e_i\le2S(A)\), so \(|e_i|<R\). Its divisibility by R forces e\_i=0, and the second difference of y is then zero as well. Conversely, evaluating two source progressions clearly gives a target progression. The pointwise inverse of E gives the inverse on tuples. If the target common difference is zero, its first two points have the same E-preimage, so both source differences are zero. The reverse implication follows by evaluation.

Finally R(A star B)=2S(A)+2R(A)S(B)+1=R(A)R(B). Both bracketings of a triple product are the same set \(A\cup R(A)B\cup R(A)R(B)C\), and the empty set has R=1. QED.

\textbf{Concrete contrapositive.} A nonconstant k-AP in H(A star B), decoded by the displayed division map, gives a nonconstant k-AP in at least one of H(A),H(B). Thus a failed product construction comes with a factor witness, rather than an unexplained failure of a heuristic.

The same product preserves the ternary image exactly. Put \(T(A)=\{\sum_a x_a a:x_a\in\{0,1,2\}\}\). Since 0\textless=x\textless=2S(A)=R(A)-1 for x in T(A), division by R gives a bijection \(T(A)\times T(B)\to T(A\star B)\). Consequently the unweighted image polynomials satisfy

\[
\begin{adjustbox}{max width=.92\linewidth}
$\displaystyle \sum_{t\in T(A\star B)}z^t=
\left(\sum_{x\in T(A)}z^x\right)
\left(\sum_{y\in T(B)}z^{R(A)y}\right).$
\end{adjustbox}
\]

This supplies a direct connection to the earlier variance work while keeping its finite sets and their measures explicit.

\section{3. Existence of every fixed-k exponential rate}\label{n03-existence-of-every-fixed-k-exponential-rate}

\textbf{Theorem 3.1.} For every integer k\textgreater=3 the following limit exists:

\[
\begin{adjustbox}{max width=.92\linewidth}
$\displaystyle \boxed{\lambda_k=\lim_{n\to\infty}g_k(n)^{1/n}
=\inf_{n\ge1}h_k(n)^{1/n}
=\inf_{\substack{A\in\mathcal A_k(|A|)\\ |A|>0}}
(2S(A)+1)^{1/|A|}.}$
\end{adjustbox}
\]

Also \(1<\lambda_k\le3\), and the sequence lambda\_k is nonincreasing in k.

\textbf{Proof.} Powers of 3 provide admissible sets for every n: their subset sums use only base-3 digits 0 and 1, and reduction of a second-difference relation modulo 3 followed by division forces each digit triple to be constant. A k-AP would have a nonconstant initial 3-AP. Thus the displayed minima exist, and h\_k(0)=1.

Theorem 2.1 gives h\_k(n+m)\textless=h\_k(n)h\_k(m). Put b\_n=log h\_k(n) and L=inf\_\{l\textgreater=1\} b\_l/l. For fixed l write n=ql+r, 0\textless=r\textless l. Repeated composition gives the exact inequality b\_n\textless=q b\_l+b\_r. Division by n and passage to the upper limit gives limsup b\_n/n\textless=b\_l/l, since the finitely many b\_r are bounded. Taking the infimum over l gives limsup\textless=L. The definition of L gives b\_n/n\textgreater=L for every n\textgreater=1. Therefore b\_n/n tends to L.

The original maximum and the sum observable obey

\[
\begin{adjustbox}{max width=.92\linewidth}
$\displaystyle 2g_k(n)+1\le h_k(n)\le2n g_k(n)+1.$
\end{adjustbox}
\]

The left inequality follows from S(A)\textgreater=max A\textgreater=g\_k(n). For the right one, use a set attaining g\_k(n). In particular

\[
\begin{adjustbox}{max width=.92\linewidth}
$\displaystyle 0\le\frac{\log h_k(n)-\log g_k(n)}n
\le\frac{\log(2n+1)}n\longrightarrow0.$
\end{adjustbox}
\]

Thus log g\_k(n)/n tends to L as well. Exponentiating proves all three expressions. The elementary construction gives lambda\_k\textless=3. Inclusion \(\mathcal A_k(n)\subseteq\mathcal A_{k+1}(n)\) follows by the explicit prefix map \(\operatorname{AP}_{k+1}(X)\to\operatorname{AP}_k(X)\), which preserves a nonzero common difference. Hence lambda is nonincreasing. Strict positivity above 1 follows from the next calculation. QED.

For completeness, the known elementary lower mechanism can be restated without an unproved input. Fix an ordering of an admissible A and let B be the partial subset-sum set just before a positive generator a is added. Split B into maximal chains with step a. A chain of length k-1 would extend in B union (B+a) to a k-AP, so all chain lengths are at most k-2. Each chain adds one point. Consequently

\[
\begin{adjustbox}{max width=.92\linewidth}
$\displaystyle |B\cup(B+a)|\ge\frac{k-1}{k-2}|B|.$
\end{adjustbox}
\]

Iteration from \{0\} and \textbar H(A)\textbar\textless=S(A)+1\textless=n max A+1 gives

\[
\begin{adjustbox}{max width=.92\linewidth}
$\displaystyle g_k(n)\ge\frac{((k-1)/(k-2))^n-1}{n},\qquad
\lambda_k\ge\frac{k-1}{k-2}>1.$
\end{adjustbox}
\]

This chain mechanism and stronger adaptive finite estimates are due to the existing literature, in particular Korsky (2026, Section 5). It is restated here only to close the elementary existence theorem\textquotesingle s bounds.

\section{4. The exact finite carry graph}\label{n03-the-exact-finite-carry-graph}

Let b\textgreater=2 and let D be a finite nonempty subset of \{0,...,b-1\}. Define

\[
\begin{adjustbox}{max width=.92\linewidth}
$\displaystyle L_m(D;b)=\{\sum_{j=0}^{m-1}d_jb^j:d_j\in D\}.$
\end{adjustbox}
\]

Words are written from least significant to most significant digit. The map \(e_m:D^m\to L_m(D;b)\) is a bijection: its inverse is the m-digit base-b expansion. No digit identification occurs in this map.

For k\textgreater=3 define the directed, labeled graph with vertices

\[
\begin{adjustbox}{max width=.92\linewidth}
$\displaystyle (c,f)\in\{-1,0,1\}^{k-2}\times\{0,1\}.$
\end{adjustbox}
\]

For a digit column \(d=(d_0,\ldots,d_{k-1})\in D^k\), let \(\Delta_i(d)=d_{i+2}-2d_{i+1}+d_i\). The column labels an edge

\[
\begin{adjustbox}{max width=.92\linewidth}
$\displaystyle (c,f)\xrightarrow{d}(c',f')$
\end{adjustbox}
\]

precisely when, for every i,

\[
\begin{adjustbox}{max width=.92\linewidth}
$\displaystyle \boxed{b c'_i=c_i+\Delta_i(d),\qquad
f'=f\mathbin{\lor}[d_0\ne d_1].}$
\end{adjustbox}
\]

The initial vertex is (0,0), and the accepting vertex is (0,1).

\textbf{Theorem 4.1 (all-length decision and witness bound).} For every b,D,k as above, the following are equivalent:

\begin{enumerate}
\def\labelenumi{(\alph{enumi})}
\item
  Some L\_m(D;b), m\textgreater=1, contains a nonconstant k-AP.
\item
  The accepting vertex is reachable from the initial vertex.
\end{enumerate}

Every such instance has a witness using at most \(3^{k-2}\) digit columns. A finite vertex set containing the initial vertex, excluding the accepting vertex, and closed under every labeled transition certifies k-AP-freeness for every m. When 0 belongs to D, it also certifies the union over all m, since leading zero padding places any finite tuple in one common L\_m.

\textbf{Proof.} A matrix of digit columns \(d^{(0)},\ldots,d^{(m-1)}\) determines numbers \(y_i=\sum_{j<m}b^j d_i^{(j)}\). Starting at c\^{}(0)=0, the edge equations say \(\Delta(d^{(j)})=b c^{(j+1)}-c^{(j)}\). Hence, exactly,

\[
\begin{adjustbox}{max width=.92\linewidth}
$\displaystyle \sum_{j<m}b^j\Delta(d^{(j)})=b^m c^{(m)}.$
\end{adjustbox}
\]

Ending with c\^{}(m)=0 is equivalent to all second differences of y vanishing. The flag is 1 exactly when the first two digit words differ; uniqueness of base-b expansion makes this equivalent to y\_0!=y\_1. Thus any accepting labeled path evaluates to a nonconstant k-AP.

Conversely, take a k-AP in L\_m and use the unique base-b expansion of each term. Set

\[
\begin{adjustbox}{max width=.92\linewidth}
$\displaystyle c^{(j)}=b^{-j}\sum_{t<j}b^t\Delta(d^{(t)}).$
\end{adjustbox}
\]

The full second difference is zero, so the prefix sum is divisible by b\^{}j; each c\^{}(j) is an integer vector. Since \textbar Delta\_i\textbar\textless=2(b-1),

\[
\begin{adjustbox}{max width=.92\linewidth}
$\displaystyle |c_i^{(j)}|\le2(1-b^{-j})<2\quad(j>0).$
\end{adjustbox}
\]

It lies in \{-1,0,1\}. These carries satisfy the transition equations and end at zero. The nonconstant progression makes the flag 1. This is the inverse map from progressions, through their unique digits, to labeled paths.

The only reachable flag-0 vertex is (0,0). Indeed, from zero with d\_0=d\_1, the congruence for d\_2 forces d\_2=d\_1 because both are in {[}0,b); successive congruences force all entries equal, and then every carry stays zero. The graph therefore has at most one reachable flag-0 vertex and 3\^{}(k-2) reachable flag-1 vertices. A shortest accepting path repeats no vertex, so its length is at most 3\^{}(k-2). Closure of a proposed safe vertex set keeps every finite path inside that set, proving the last assertions. QED.

The graph may be explored using only \textbar D\textbar\^{}2 starting digit pairs at each vertex. Having chosen d\_0,d\_1, each subsequent digit is forced by

\[
\begin{adjustbox}{max width=.92\linewidth}
$\displaystyle d_{i+2}\equiv 2d_{i+1}-d_i-c_i\pmod b,
\qquad 0\le d_{i+2}<b.$
\end{adjustbox}
\]

An edge exists exactly when all these forced digits belong to D. This is an exact enumeration of D\^{}k columns satisfying the divisibility conditions. The checker also compares it with direct column enumeration on a smaller, fully specified domain. Parent edges retain actual digit columns so that negative results contain concrete integer witnesses.

\subsection{4.2 Retaining digits above the base and every representation collision}\label{n03-retaining-digits-above-the-base-and-every-representation-collision}

The preceding canonical-digit map has a complete extension to arbitrary finite D contained in {[}0,S{]}, with S allowed to exceed b. Keep the same evaluation surjection e\_m:D\^{}m -\textgreater{} L\_m(D;b), and retain its entire fibers. For example, when D=\{0,7,8,15,19,26,27,34\} and b=19, the words (19,7) and (0,8) both belong to D\^{}2 and evaluate to 152. Both are retained. Tracking whether two digit words differ would flag these equal values incorrectly. The map to actual differences is therefore made explicit below.

Set

\[
\begin{adjustbox}{max width=.92\linewidth}
$\displaystyle C_0=\left\lceil\frac{2S}{b-1}\right\rceil,\qquad
H_0=\left\lceil\frac{S}{b-1}\right\rceil.$
\end{adjustbox}
\]

The full graph has vertices

\[
\begin{adjustbox}{max width=.92\linewidth}
$\displaystyle (c,h,f)\in\{-C_0,\ldots,C_0\}^{k-2}
\times\{-H_0,\ldots,H_0\}\times\{0,1\}.$
\end{adjustbox}
\]

For a column d in D\^{}k require the exact transition equations

\[
\begin{adjustbox}{max width=.92\linewidth}
$\displaystyle bc'_i=c_i+\Delta_i(d),\quad
h+d_1-d_0=r+bh',\quad 0\le r<b,\quad
f'=f\lor[r\ne0].$
\end{adjustbox}
\]

Thus h\textquotesingle{} and r are respectively the Euclidean quotient and remainder of h+d\_1-d\_0; negative values are handled by the floor quotient. Start at (0,0,0). A vertex is accepting exactly when c=0 and either f=1 or h!=0.

\textbf{Theorem 4.2 (full carry equivalence).} For every b\textgreater=2, k\textgreater=3, and nonempty finite D contained in the nonnegative integers, an L\_m(D;b) contains a nonconstant k-AP exactly when this full graph has a reachable accepting vertex. A shortest witness has at most

\[
\begin{adjustbox}{max width=.92\linewidth}
$\displaystyle 2(2C_0+1)^{k-2}(2H_0+1)-1$
\end{adjustbox}
\]

columns. A finite transition-closed set of vertices containing the initial vertex and excluding every accepting vertex certifies all word lengths.

\textbf{Proof.} For a labeled path of length m let \(y_i=\sum_{j<m}b^j d_i^{(j)}\). The second-difference transitions telescope exactly as before to \(\Delta(y)=b^m c^{(m)}\). Meanwhile the extra carry transitions give the separate identity

\[
\begin{adjustbox}{max width=.92\linewidth}
$\displaystyle \boxed{y_1-y_0=\sum_{j<m}b^j r_j+b^m h^{(m)}.}$
\end{adjustbox}
\]

The sum of remainders lies in {[}0,b\^{}m). Consequently this difference vanishes exactly when all r\_j vanish and h\^{}(m)=0. This is precisely the complement of the stated acceptance test at c\^{}(m)=0. Thus every accepting path gives an actual nonconstant progression, with no assumption about uniqueness of word representations.

Conversely, take any progression in L\_m and any digit word representing each of its terms. For each j define the second-difference carry from the prefix formula in Theorem 4.1. It is integral because the full second difference vanishes. Its magnitude is at most \(2S(1-b^{-j})/(b-1)\), hence at most C\_0. Starting with h\^{}(0)=0, perform the displayed quotient-remainder recursion. If \(p_j=\sum_{t<j}b^t(d_1^{(t)}-d_0^{(t)})\), induction gives \(h^{(j)}=\lfloor p_j/b^j\rfloor\) and a remainder in {[}0,b\^{}j). Since \(|p_j|/b^j\le S(1-b^{-j})/(b-1)\), one has \(-H_0\le h^{(j)}\le H_0\). These data define a full-graph path. The actual nonzero difference guarantees its accepting endpoint.

This construction is a bijection between labeled accepting paths and matrices of digit words whose evaluated rows form nonconstant progressions. Evaluation of the rows maps these matrices surjectively to the progression set, with exactly the original representation fibers. A shortest accepting path has no repeated vertex, proving the finite bound. Closure proves the all-length assertion. QED.

For canonical D, forgetting h gives a graph morphism on reachable states from this graph to the graph of Theorem 4.1: (c,h,f) -\textgreater{} (c,f). Before the first differing canonical digit, h=0; at the first such digit the difference has magnitude less than b and gives a nonzero remainder. Hence both flags become true on exactly the same column. This proves preservation of the transition labels and the acceptance condition on reachable vertices.

An explicitly checked positive full certificate uses b=19, A=\{7,8,19\}, and D=H(A)=\{0,7,8,15,19,26,27,34\}. Its full graph has 15 reachable states and 288 labeled transitions. Direct D\^{}4 enumeration verifies closure. There is also an exact independent generator-set explanation:

\[
\begin{adjustbox}{max width=.92\linewidth}
$\displaystyle \bigcup_{j<m}19^j\{7,8,19\}
\ \subseteq\ \bigcup_{j\le m}19^j\{1,7,8\}.$
\end{adjustbox}
\]

The terms 7 and 8 stay at their original levels, while each 19 at level j is the original 1-generator at level j+1. The subset-sum inclusion sends any progression witness to one in the original certified construction. This example uses digits above the base without replacing the input block.

\section{5. Modular reduction, carry compatibility, and explicit contrapositives}\label{n03-modular-reduction-carry-compatibility-and-explicit-contrapositives}

Reduction \(\rho_b:\mathbb Z\to\mathbb Z/b\mathbb Z\) is a group homomorphism. Applied coordinatewise to a digit column, it maps zero-start carry transitions to modular arithmetic progressions, since \(\Delta(d)=bc'\). The target carry records the exact integer second difference divided by b. A sequence of such columns lifts to an integer progression precisely when the carries match and eventually return to zero, as proved by Theorem 4.1.

In particular, a modularly k-AP-free D has only constant zero-start columns; \{(0,0)\} is then a closed safe certificate. This proves the usual modular digit lifting argument with no assumption that b is prime.

Two explicit examples retain the additional information carried by the graph. For b=4, D=\{0,2\}, k=3, the complete reachable set is

\[
\begin{adjustbox}{max width=.92\linewidth}
$\displaystyle \{(0,0),(-1,1),(1,1)\}.$
\end{adjustbox}
\]

The modular tuple (0,2,0) labels an edge from (0,0) to (-1,1). The vertex (-1,1) has no outgoing edge: every integer second difference of D-digits is even, so adding -1 cannot be divisible by 4. The same holds at (1,1). Consequently the whole digit language avoids 3-APs even though that modular tuple exists. This uses the original digits \{0,2\} exactly as given.

For b=8, D=\{0,1,3,4\}, k=5, the complete reachable set is

\[
\begin{adjustbox}{max width=.92\linewidth}
$\displaystyle \{((0,0,0),0),((-1,1,-1),1),((1,-1,1),1)\}.$
\end{adjustbox}
\]

Only constant columns occur as zero-to-zero transitions. The other two reachable vertices have no outgoing edge; their closure can be checked by the 16 possible first digit pairs at each vertex. The reduction map sends (0,4,0,4,0) to a modular progression, and the exact carry (-1,1,-1) records its terminal incompatibility. Thus arbitrary nonzero carries are retained rather than excluded by assumption.

The opposite outcome is equally explicit. For the four-generator block \{1,3,4,7\} in base 16, the digit columns

\[
\begin{adjustbox}{max width=.92\linewidth}
$\displaystyle (0,4,8,12,0),\qquad(0,0,0,0,1)$
\end{adjustbox}
\]

form an accepting path and evaluate to (0,4,8,12,16). For the original \{1,7,8\} block in base 18, columns (0,9,0,9),(0,0,1,1) evaluate to (0,9,18,27). These are witnesses against those exact smaller-base lifts. They leave the valid base-23 and base-19 constructions intact.

\section{6. A capacity formula covering every admissible generator set}\label{n03-a-capacity-formula-covering-every-admissible-generator-set}

A \textbf{modular certificate} is a pair (b,A) with b\textgreater=2, nonempty A a finite set of distinct positive integers, S(A)\textless b, and H(A) modularly k-AP-free. A \textbf{carry certificate} has the same data, replacing the last property by absence of an accepting path in the graph of Section 4 for D=H(A). Call their classes \(\mathcal M_k\) and \(\mathcal C_k\).

\textbf{Theorem 6.1 (exact certificate capacity).} For every k\textgreater=3,

\[
\begin{adjustbox}{max width=.92\linewidth}
$\displaystyle \boxed{\lambda_k
=\inf_{(b,A)\in\mathcal M_k}b^{1/|A|}
=\inf_{(b,A)\in\mathcal C_k}b^{1/|A|}.}$
\end{adjustbox}
\]

Both classes admit finite exact verification of each member.

\textbf{Proof.} From either certificate form \(A^{[m]}=\bigcup_{j=0}^{m-1}b^jA\). Its members are distinct: each a in A is in {[}1,b), so two members at different levels have disjoint magnitude ranges. Thus \(|A^{[m]}|=m|A|\) and \(\max A^{[m]}=b^{m-1}\max A\).

The blockwise subset-sum map \(\{0,1\}^{m\times |A|}\to H(A)^m\) is surjective, with fiber cardinality \(\prod_j\mu_A(d_j)\). Composing with e\_m gives exactly \(H(A^{[m]})=L_m(H(A);b)\). The actual subset collisions therefore persist with known multiplicities, while the value-word map is bijective. Section 4 or 5 proves k-AP-freeness at every m. Deleting generators gives, for every n\textgreater=1,

\[
\begin{adjustbox}{max width=.92\linewidth}
$\displaystyle \boxed{g_k(n)\le(\max A)\,b^{\lceil n/|A|\rceil-1}.}$
\end{adjustbox}
\]

Taking roots and using Theorem 3.1 proves lambda\_k\textless=b\^{}(1/\textbar A\textbar), for either certificate class.

Conversely, for every nonempty admissible A, take b=2S(A)+1. Any modular k-AP in the integer digit set H(A) has second differences divisible by b and of absolute value at most 2S(A)=b-1. They therefore vanish over the integers. Admissibility makes the progression constant. Thus A -\textgreater{} (2S(A)+1,A) is an explicit map from every admissible set to a modular certificate. Theorem 3.1 makes the infimum of the resulting costs exactly lambda\_k. Since every modular certificate is also a carry certificate, the opposite inequalities are proved. QED.

\textbf{Finite countercertificate equivalence.} For every positive real u,

\[
\begin{adjustbox}{max width=.92\linewidth}
$\displaystyle \lambda_k<u\quad\Longleftrightarrow\quad
\exists(b,A)\in\mathcal M_k\;[b<u^{|A|}].$
\end{adjustbox}
\]

The same equivalence holds with C\_k. This follows from the two infimum identities, including the strict inequality on both sides. It is a proved witness statement, without an unproved extension premise. For an algebraic threshold u=c\^{}(1/t), positive integers c,t, the comparison is the exact integer inequality \(b^t<c^{|A|}\).

There is a concrete monotone sequence approaching lambda\_k from above: minimize b\^{}(1/\textbar A\textbar) over certificates with 2\textless=b\textless=B. For each B this domain is finite, because A is a subset of {[}1,b-1{]} with S(A)\textless b and \textbar A\textbar(\textbar A\textbar+1)/2\textless b. For B\textgreater=3 it is nonempty. Exhausting B exhausts all certificates, so its limit equals lambda\_k. The theorem supplies convergence; this contribution supplies no effective error bound for that convergence.

\subsection{6.2 Full certificates and exact generator distinctness}\label{n03-full-certificates-and-exact-generator-distinctness}

A finite positive set A is \textbf{b-separated} when no two of its elements have ratio b\^{}t for an integer t\textgreater=1. This is a finite decidable predicate: for each a, multiply successively by b until exceeding max A, checking membership at each step. The map

\[
\begin{adjustbox}{max width=.92\linewidth}
$\displaystyle q_m:\{0,\ldots,m-1\}\times A\longrightarrow
\bigcup_{j<m}b^jA,\qquad q_m(j,a)=b^ja$
\end{adjustbox}
\]

is surjective. It is injective for every m exactly when A is b-separated. Indeed, equality b\^{}i a=b\^{}j a\textquotesingle{} either gives i=j and a=a\textquotesingle, or, after taking the smaller index, gives precisely a power ratio. Conversely, a\textquotesingle=b\^{}t a gives the explicit collision q\_(t+1)(0,a\textquotesingle)=q\_(t+1)(t,a).

A \textbf{full certificate} consists of b\textgreater=2 and a nonempty b-separated set A of distinct positive integers, together with a safe full graph for D=H(A). There is no restriction S(A)\textless b in this definition. Write its class as F\_k.

\textbf{Theorem 6.2.} For every k\textgreater=3,

\[
\begin{adjustbox}{max width=.92\linewidth}
$\displaystyle \boxed{\lambda_k=\inf_{(b,A)\in\mathcal F_k}b^{1/|A|}.}$
\end{adjustbox}
\]

Every full certificate gives, for all n\textgreater=1, \(g_k(n)\le(\max A)b^{\lceil n/|A|\rceil-1}\).

\textbf{Proof.} Injectivity of q\_m gives exactly m\textbar A\textbar{} distinct generators. The map from subsets of these generators to their block digit words is surjective, and evaluation gives their actual subset-sum image. Because word evaluation may now have multiple fibers, the full multiplicity identity is

\[
\begin{adjustbox}{max width=.92\linewidth}
$\displaystyle \mu_{A^{[m]}}(y)=
\sum_{\substack{d\in H(A)^m\\\sum_j b^j d_j=y}}
\prod_{j<m}\mu_A(d_j).$
\end{adjustbox}
\]

This follows directly by partitioning the binary generator choices according to their block sums. Theorem 4.2 certifies the actual value image. The maximum is (max A)b\^{}(m-1), and deletion gives the stated finite bound. Therefore lambda\_k\textless=b\^{}(1/\textbar A\textbar) for each full certificate. Every canonical carry certificate is a full certificate, because its generators lie in {[}1,b) and its digit language has the same progression property under the full equivalence. Theorem 6.1 provides the reverse infimum inequality. QED.

The separation check is retained even when a digit-language check is available. For b=19 and A=\{1,7,8,19\}, q\_m has m-1 duplicate pairs, and its image has 3m+1 elements, compared with 4m indexed generator positions. For instance, q\_2(0,19)=q\_2(1,1)=19. The actual image embeds in the indexed system by choosing, for each image value, its least-level preimage; the resulting injection of finite subsets preserves their sums. It need not cover arbitrary indexed selections. The verifier detects this power alias before accepting a generator-cardinality certificate; it never silently deletes a generator and retains the original cardinality claim.

\section{7. Explicit applications beyond k=4}\label{n03-explicit-applications-beyond-k4}

The small certificates are

\[
\begin{adjustbox}{max width=.92\linewidth}
$\displaystyle (k,b,A)=(5,23,\{1,3,4,7\}),\qquad
(6,47,\{1,2,6,7,14\}).$
\end{adjustbox}
\]

Their digit sets are respectively

\[
\begin{adjustbox}{max width=.92\linewidth}
$\displaystyle D_{23}=\{0,1,3,4,5,7,8,10,11,12,14,15\}$
\end{adjustbox}
\]

and

\[
\begin{adjustbox}{max width=.92\linewidth}
$\displaystyle D_{47}=\{0,1,2,3,6,7,8,9,10,13,14,15,16,17,20,21,22,23,24,27,28,29,30\}.$
\end{adjustbox}
\]

To make the finite modular fact inspectable, for each positive step up to half the prime modulus the maximum number of consecutive terms in D is:

\begin{longtable}[]{@{}lll@{}}
\toprule\noalign{}
b & steps d in order & maximum run lengths in that order \\
\midrule\noalign{}
\endhead
\bottomrule\noalign{}
\endlastfoot
23 & 1,...,11 & 3,4,4,4,4,2,3,4,3,4,4 \\
47 & 1,...,23 & 5,3,4,4,3,5,5,4,3,4,3,3,3,3,3,2,4,2,4,4,4,3,4 \\
\end{longtable}

Each row is verified by starting at every residue and iterating the stated step until leaving D. The reversal bijection \((x_0,...,x_{k-1})\mapsto(x_{k-1},...,x_0)\) exchanges d and -d, so the table covers every nonzero step. The executable additionally checks all b(b-1) ordered start/step pairs directly, without using that reduction.

A larger exact search gives the six-generator block

\[
\begin{adjustbox}{max width=.92\linewidth}
$\displaystyle A_* = \{1,4,5,17,21,22\},\qquad S(A_*)=70.$
\end{adjustbox}
\]

Its 38 distinct subset sums are

\[
\begin{adjustbox}{max width=.92\linewidth}
$\displaystyle \begin{aligned}
D_* =\{&0,1,4,5,6,9,10,17,18,21,22,23,25,26,27,28,30,31,32,\\
       &38,39,40,42,43,44,45,47,48,49,52,53,60,61,64,65,66,69,70\}.
\end{aligned}$
\end{adjustbox}
\]

This is modularly 5-AP-free at b=97 and modularly 6-AP-free at b=93. The second check includes all nonzero steps of the composite modulus, including steps of order 3 and 31. In each case the complete reachable carry certificate consists of the initial vertex alone. The exhaustive finite arithmetic is included in the checker and receipt.

\textbf{Theorem 7.1 (new explicit upper constructions).} For every n\textgreater=1,

\[
\begin{adjustbox}{max width=.92\linewidth}
$\displaystyle \boxed{g_5(n)\le22\,97^{\lceil n/6\rceil-1},\qquad
       g_6(n)\le22\,93^{\lceil n/6\rceil-1}.}$
\end{adjustbox}
\]

Consequently \(\lambda_5\le97^{1/6}\) and \(\lambda_6\le93^{1/6}\). These statements follow by applying the fully proved lift to the displayed finite certificates. The smaller certificates also give \(g_5(n)\le7\cdot23^{\lceil n/4\rceil-1}\) and \(g_6(n)\le14\cdot47^{\lceil n/5\rceil-1}\), useful for some finite n. The minimum of the corresponding finite upper expressions may be used.

The target progression length is retained in the certificate. At base 93, the same block\textquotesingle s two-level subset sums contain the five-term progression

\[
\begin{adjustbox}{max width=.92\linewidth}
$\displaystyle 9,\ 990,\ 1971,\ 2952,\ 3933$
\end{adjustbox}
\]

with difference 981. Its digit columns are (9,60,18,69,27) and (0,10,21,31,42), all entries in \(D_*\). The prefix map from six-term progression tuples to five-term tuples is the restriction of Z\^{}6 -\textgreater{} Z\^{}5. It preserves the common difference. This particular five-term tuple has no preimage: its forced sixth term is 4914, whose residue 78 modulo 93 is absent from \(D_*\). Thus the base-93 result is certified for k=6, while base 97 supplies the stated k=5 construction.

For a sharper remainder-sensitive version, list a certificate\textquotesingle s generators as a\_1\textless...\textless a\_s, write n=sq+r with 0\textless=r\textless s, and retain q complete levels and then the first r generators at level q. When r\textgreater0 the maximum is a\_r b\^{}q; when r=0 it is a\_s b\^{}(q-1). Positivity and disjoint level ranges prove the cardinality and the bound exactly.

The comparisons with Korsky\textquotesingle s stated general graph bound are exact: for k=5 its rate is sqrt(5), and 97\textless5\^{}3; for k=6 its rate is 7\^{}(2/5), and \(93^5<7^{12}\). These comparisons identify improvement over that retrieved bound, without a claim of global literature priority.

\section{8. An explicit correlated cube behind the six-generator block}\label{n03-an-explicit-correlated-cube-behind-the-six-generator-block}

Let the six columns in Z\^{}3 be

\[
\begin{adjustbox}{max width=.92\linewidth}
$\displaystyle v_1=(1,0,0),\ v_2=(0,1,0),\ v_3=(1,1,0),\
v_4=(0,0,1),\ v_5=(0,1,1),\ v_6=(1,1,1).$
\end{adjustbox}
\]

Let C be the image of the binary cube under \(\varepsilon\mapsto
\sum_i\varepsilon_i v_i\). It lies in {[}0,3{]} x {[}0,4{]} x {[}0,3{]}.

\textbf{Proposition 8.1.} C has no nonconstant five-term arithmetic progression over Z\^{}3.

\textbf{Proof.} In a five-term progression, the first and third coordinates have integer common differences and span at least 4 times their absolute values. Their ranges have width 3, so both common differences are zero. A remaining nonzero middle difference must be +/-1 and use both middle levels 0 and 4. Middle level 0 forces the choices for columns 2,3,5,6 all to vanish; its first coordinate is at most 1. Middle level 4 forces all four of those choices to be present; its first coordinate is at least 2. They cannot share the required first coordinate. Thus every common difference is zero. QED.

The homomorphism \(L:\mathbb Z^3\to\mathbb Z\), L(x,y,z)=x+4y+17z, maps these columns to \(A_*\). Direct enumeration gives \textbar C\textbar=38 and \textbar L(C)\textbar=38, so its restriction C -\textgreater{} \(H(A_*)\) is a bijection. However the second-difference equations after projection require their own verification: \textbar C+C\textbar=201 whereas \textbar L(C+C)\textbar=139. The exact surjection L:C+C -\textgreater{} T(A\_*) retains this distinction; its fibers are included in the executable checks. The modular certificate in Section 7 verifies the required progression property after projection. The three-dimensional argument motivates the construction without silently asserting a Freiman equivalence that the sumset cardinalities contradict.

\section{9. Overlapping relations of arbitrary extent already occur at k=5}\label{n03-overlapping-relations-of-arbitrary-extent-already-occur-at-k5}

\textbf{Theorem 9.1.} For every r\textgreater=2 the set

\[
\begin{adjustbox}{max width=.92\linewidth}
$\displaystyle A_r=\{5^i:0\le i<r\}\cup\{5^i+5^{i+1}:0\le i<r-1\}$
\end{adjustbox}
\]

has 2r-1 distinct positive generators, has 5-AP-free subset sums, and has minimal signed relations whose intersection graph contains a path of r-1 vertices. Consequently the relation interaction has arbitrarily large connected extent for every k\textgreater=5.

\textbf{Proof.} Label generators v\_i=5\^{}i and \(e_i=5^i+5^{i+1}\). Their standard base-5 digit vectors are respectively the unit vector u\_i and u\_i+u\_(i+1), which proves distinctness. The exact column homomorphism from Z\^{}\{2r-1\} to Z\^{}r sends the generator labels to those columns; evaluation \(\eta(x)=\sum_i x_i5^i\) gives the actual integer weights. Each subset contributes at most three units at each position, one from its vertex and at most two from incident edges. Thus H(A\_r) lies in the base-5 language with digits \{0,1,2,3\}, which is modularly 5-AP-free and hence safe by Section 5.

The signed relation q\_i has coefficients +1 on v\_i,v\_(i+1), -1 on e\_i, and zero elsewhere. It is minimal: a relation supported on a proper nonempty subset would have one or two coordinates, incompatible with positivity and distinctness of the weights. Consecutive q\_i share v\_(i+1), producing the claimed path. The map \(\Gamma:\mathbb Z^{r-1}\to\ker(\Phi_{A_r})\), \(\Gamma(t)=\sum_i t_iq_i\), is injective because its e\_i coordinate is -t\_i. This is an explicit linear embedding of the entire path relation system into the actual generator kernel. Admissibility for larger k follows from the prefix map on progression tuples. QED.

The old four-term splitting algebra remains visible. For adjacent edges, q\_i+q\_(i+1) has magnitude-two layer consisting of the single nonzero weight v\_(i+1). It gives the actual four-term progression

\[
\begin{adjustbox}{max width=.92\linewidth}
$\displaystyle v_i+v_{i+2},\quad v_i+v_{i+1}+v_{i+2},\quad
e_i+e_{i+1},\quad e_i+e_{i+1}+v_{i+1},$
\end{adjustbox}
\]

with common difference v\_(i+1). At r=3 these values are 26,31,36,41 in H(\{1,5,25,6,30\}). The restriction of the coordinate projection Z\^{}5 -\textgreater{} Z\^{}4 gives the prefix map \(\operatorname{AP}_5(H(A_r))\to\operatorname{AP}_4(H(A_r))\) preserves common differences; the displayed 4-AP lies outside its image. Its witness and the 5-AP-free proof identify exactly which step of a four-term-only decomposition cannot be imposed on the general-k domain.

\section{10. A proved boundary for all positive interval-digit column models}\label{n03-a-proved-boundary-for-all-positive-interval-digit-column-models}

\textbf{Proposition 10.1 (full positive-column budget).} Let M be an r by n nonnegative integer matrix with distinct nonzero columns. Let every row sum be at most tau. Then

\[
\begin{adjustbox}{max width=.92\linewidth}
$\displaystyle \boxed{2n\le r(\tau+1).}$
\end{adjustbox}
\]

\textbf{Proof.} At most r columns have coordinate sum 1, since they must be unit vectors. Every other column has coordinate sum at least 2. The total column mass is therefore at least 2n-r and is exactly the total row mass, which is at most r tau. Combining the exact counts proves the inequality. QED.

The column map M:Z\^{}n -\textgreater{} Z\^{}r and the evaluation homomorphism \(\eta_b(x)=\sum_{i=0}^{r-1}x_i b^i\) give the integer generators. When tau\textless b, evaluation is injective on the full subset-coordinate box; its inverse is the actual base-b digit expansion. For prime b and tau=min(b,k)-2 this is the positive interval-digit construction. The budget forces r\textgreater=ceil(2n/(min(b,k)-1)). With r active positions, some generator has value at least b\^{}(r-1). Thus, for each fixed prime b, that entire column class has exponential rate at least b\^{}(2/(min(b,k)-1)); the graph construction attains that rate asymptotically. This is a statement about the precisely displayed column domain and maps.

An explicit failure of a proposed universal passage into that domain is supplied by A=\{1,3,4,7\}. For any assignment of nonzero nonnegative columns w\_1,w\_3,w\_4,w\_7 preserving its subset equalities,

\[
\begin{adjustbox}{max width=.92\linewidth}
$\displaystyle w_4=w_1+w_3,\qquad w_7=w_3+w_4,
\qquad \sum_{a\in A}w_a=3w_1+4w_3.$
\end{adjustbox}
\]

Some coordinate of w\_3 is at least 1, and that row load is at least 4. Hence every such assignment violates a row-load bound of 3. The typed relation-preserving map M:Z\^{}4 -\textgreater{} Z\^{}r has to send the two displayed relation vectors to zero; those equations prove the obstruction in every dimension. The actual scalar map (coefficients 1,3,4,7) and the modular projection to Z/23 are exhibited in Section 7 and produce a valid 5-AP-free construction. Thus the carried and correlated digit route has a concrete mathematical reason to be pursued beyond positive interval boxes.

\section{11. What has been established and what remains to be evaluated}\label{n03-what-has-been-established-and-what-remains-to-be-evaluated}

The proved all-k conclusions are the exact product, the existence of each lambda\_k, the complete certificate-capacity formulas, and the finite canonical and full carry decisions with an explicit witness-length bound. The explicit k=5 and k=6 certificates improve the retrieved general upper estimates. The path construction and positive-column count prove two precise structural boundaries for possible lower-bound arguments.

The exact numerical values of lambda\_k for general k\textgreater=5, and the sharp large-k behavior of log lambda\_k, remain unevaluated here. The current lower and upper mechanisms leave a gap between order 1/k and order log(k)/k. The new capacity identity identifies the exact optimization across all finite certificates, including every admissible integer set through its explicit map A -\textgreater{} (2S(A)+1,A). A bounded search supplies exact finite data and decreasing upper bounds; it supplies no lower certificate for unsearched bases. Independent review should focus on Sections 2,4,6 before using the capacity theorem, and on the modular transcripts before citing new constants.

\section{12. Reproducible verification and exact search domains}\label{n03-reproducible-verification-and-exact-search-domains}

The standalone command, run from this contribution\textquotesingle s directory, is

\begin{Verbatim}[breaklines=true,breakanywhere=true,fontsize=\footnotesize]
python certificates/verify_general_k.py --output certificates/general_k_receipt.json
\end{Verbatim}

The resulting receipt reports PASS for the following completely specified finite domains. The general implications are proved in Sections 2-\/-6; these checks independently test their finite maps, graph transitions, and witnesses.

\begin{longtable}[]{@{}
  >{\raggedright\arraybackslash}p{(\columnwidth - 2\tabcolsep) * \real{0.5000}}
  >{\raggedright\arraybackslash}p{(\columnwidth - 2\tabcolsep) * \real{0.5000}}@{}}
\toprule\noalign{}
\begin{minipage}[b]{\linewidth}\raggedright
Check
\end{minipage} & \begin{minipage}[b]{\linewidth}\raggedright
Domain and result
\end{minipage} \\
\midrule\noalign{}
\endhead
\bottomrule\noalign{}
\endlastfoot
Primary modular certificates & Every start and nonzero step: 342 pairs for base 19; 506 for base 23; 2,162 for base 47; 9,312 for base 97; 8,556 for base 93 \\
Canonical carry graph & All D subset {[}0,b) containing zero, 2\textless=b\textless=8, 3\textless=k\textless=6: 1,016 instances; 415 safe and 601 with explicit witnesses \\
Independent canonical closure & Direct enumeration of digit columns verifies 1,433 transitions across the safe certificates \\
Full carry graph & All D subset {[}0,6{]} containing zero, b=2,3,4, k=3,4: 384 instances; 30 safe and 354 with explicit witnesses \\
Independent full closure & 122 transitions across the small safe cases, plus 288 transitions for the noncanonical \{7,8,19\} example \\
Direct finite languages & All graph-domain cases at word lengths 1,2,3: 4,200 comparisons with separately generated integer sets \\
Exact composition & 256 ordered pairs of empty, singleton, or two-element subsets of {[}1,5{]}; 1,024 full AP-tuple comparisons for k=3,4,5,6 \\
Associativity & All 4,096 ordered triples from the same 16-set domain \\
Primary generator lifts & Two levels for each of five primary certificates; complete binary fiber multiplicities and remainder-sensitive deletions \\
Structural maps & The 38-point correlated cube, all 201 points in its sumset and all 139 projected values; path-relation family r=2,3,4,5 \\
Regressions & Two modular/carry bridge examples, five decoded failed-base/progression-length constructions, power-alias cardinalities, and nine rejected invalid inputs or mutations \\
\end{longtable}

The canonical and full graph outputs are additionally compared on every canonical input with b\textless=6 and k\textless=5. Two independent transition generators are used: the search solver builds columns recursively by congruences, while the closure checker directly enumerates D\^{}k. Negative outputs retain digit columns, which the checker evaluates independently to the displayed integer progression.

The constructive search is exhaustive over every positive generator set with sum strictly less than its base for k=5,6 and bases 3 through 160. A second search uses modular certificates for k=7 through 12 and bases 3 through 80. These are 784 base/length parameter records. The generic search has a depth cap of eight, but its largest returned generator counts are six in the first domain and seven in the second. The cap therefore leaves no larger admissible set hidden: any such set would have an admissible eight- element subset. The verifier separately rechecks all 784 recorded witnesses. It does not replay the entire maximality search by default.

The exhaustive-search pruning has explicit justifications. Generators are visited in strictly increasing order. For a partial tuple, the minimum sum needed to exceed the current best cardinality is an arithmetic-series sum; branches exceeding the base budget are rejected. A failed partial digit set cannot become safe after adding generators, because its progression word remains in every larger alphabet. The search retains the actual generators and only caches the progression property of an already computed digit set. No multiplication by units, gcd removal, or re-centering of an input set is used. All comparisons of exponential costs are exact cross-power integer comparisons.

The search can be replayed with

\begin{Verbatim}[breaklines=true,breakanywhere=true,fontsize=\footnotesize]
python search_carry.py --mode carry --min-base 3 --max-base 80 --k 5 6
python search_carry.py --mode carry --min-base 81 --max-base 160 --k 5 6
python search_carry.py --mode modular --min-base 3 --max-base 80 --k 7 8 9 10 11 12
\end{Verbatim}

A further retained full-carry probe fixes A\_* and tests bases 23 through 70 for both k=5 and k=6. All 96 instances produce two-level progression witnesses. The complete witnesses are in \nolinkurl{logs/noncanonical_fixed_block.json}; this is a failure of those precise base/block pairs, with the map to integer progressions supplied by Theorem 4.2. It provides no obstruction to different blocks at the same bases.

A repeated verifier run with Python\textquotesingle s optimization flag \texttt{-O} gives the same receipt bytes. The program uses explicit exception-raising requirements, so optimization cannot remove its checks. Source and transcript SHA-256 hashes identify exact inputs and outputs. Lean elaboration was not run in this contribution; its certificate type is an ordinary mathematical proof plus finite exact verification, and existing Lean artifacts are preserved.

\section{References, inspected sources, and contribution lineage}\label{n03-references-inspected-sources-and-contribution-lineage}

Costa, S. (2026). \emph{A negative answer to the Erdős--Sárkőzy question} (Version 1) {[}Preprint{]}. arXiv:2609.06303. Source: \url{https://arxiv.org/html/2609.06303v1}. Inspected HTML, especially abstract and Section 1; used solely for the current historical k=3 question.

Korsky, S. (2026). \emph{Arithmetic progression-free subset-sum sets} (Version 1) {[}Preprint{]}. arXiv:2606.24139. Source: \url{https://arxiv.org/html/2606.24139v1}. Inspected HTML, Sections 1,5,6; used for the problem definition, known chain mechanism, and the comparison upper bound.

Anonymous/deleted Reddit contributor, \& The Clankers. (2026, September 13). \emph{The exponential rate for four-term-progression-free subset-sum sets} {[}Workbench manuscript{]}. \texttt{paper/ep817\_k4\_rate.tex} at the base revision above. The original construction and proof route belong to the anonymous contributor.

sneed-and-feed. (2026, September 13). \emph{Formalize finite upper bound theorem and core helper lemmas in Extended.lean} {[}Pull request \#1{]}. Same workbench.

The Clankers. (2026, September 14). \emph{An unweighted ternary-variance refinement of the k=4 lower bound} {[}Draft pull request \#2{]}, commit \nolinkurl{612020848e4bbec5f5a723a64c5fdc85636b5e83}. Parallel predecessor contribution; its separate review and Lean status are retained.

PolyClank methodology: \nolinkurl{erdos-straus-foundation/POLYCLANK.md} (inspected) and \nolinkurl{yang-mills-interacting-workbench/docs/polyclank/RESEARCH_STATE.md} at \nolinkurl{7ba7ea5a5dab84d88934477de93dda71f6af0662} (inspected). These supply the public research-record convention, with claim-local evidence and source lineage; they supply no additional mathematical theorem used in this note.
